% refs.tex
% ==============================================================================
% REFERÊNCIAS BIBLIOGRÁFICAS
% ==============================================================================
\section*{Referências}
\addcontentsline{toc}{section}{Referências}

\begin{itemize}[leftmargin=*]
    \item KUROSE, James F.; ROSS, Keith W. \textit{Redes de Computadores e a Internet: Uma Abordagem Top-Down}. 8. ed. São Paulo: Pearson Education do Brasil, 2021.

    \item TANENBAUM, Andrew S.; WETHERALL, David J. \textit{Redes de Computadores}. 5. ed. São Paulo: Pearson Prentice Hall, 2011.

    \item STEVENS, W. Richard. \textit{TCP/IP Illustrated, Volume 1: The Protocols}. Reading: Addison-Wesley, 1994.

    \item STEVENS, W. Richard; FENNER, Bill; RUDOFF, Andrew M. \textit{UNIX Network Programming, Volume 1: The Sockets Networking API}. 3. ed. Boston: Addison-Wesley, 2004.

    \item PYTHON SOFTWARE FOUNDATION. \textit{socket --- Low-level networking interface}. Python 3 Standard Library Documentation. Disponível em: \url{https://docs.python.org/3/library/socket.html}. Acesso em: nov. 2026.

    \item VENÂNCIO, Giovanni. \textit{Especificação do Trabalho Prático: Sistema de Balanceamento de Cargas}. Disciplina de Redes de Computadores II (CI1061), Departamento de Informática, Universidade Federal do Paraná (UFPR), 2026/2.
\end{itemize}
